\chapter{Исходные данные для конкретных расчетов}

В соответствии с индивидуальным заданием для варианта №~11 рассматривается система трех линейных алгебраических уравнений с тремя неизвестными ($n = 3$):
\begin{equation}
    \begin{cases}
        2{,}7x_1 + 0{,}9x_2 - 1{,}5x_3 = 3{,}5; \\
        4{,}5x_1 - 2{,}8x_2 + 6{,}7x_3 = 2{,}6; \\
        5{,}1x_1 + 3{,}7x_2 - 1{,}4x_3 = -0{,}14.
    \end{cases}
    \label{eq:var11}
\end{equation}

Матрица коэффициентов системы $A$ и вектор свободных членов $b$ имеют вид:
\begin{equation}
    A = \begin{pmatrix}
        2{,}7 & 0{,}9 & -1{,}5 \\
        4{,}5 & -2{,}8 & 6{,}7 \\
        5{,}1 & 3{,}7 & -1{,}4
    \end{pmatrix},
    \qquad
    b = \begin{pmatrix}
        3{,}5 \\
        2{,}6 \\
        -0{,}14
    \end{pmatrix}.
\end{equation}

Расширенная матрица системы $\widetilde{A} = (A\,|\,b)$ записывается в виде:
\begin{equation}
    \widetilde{A} = \begin{pmatrix}
        2{,}7 & 0{,}9 & -1{,}5 & | & 3{,}5 \\
        4{,}5 & -2{,}8 & 6{,}7 & | & 2{,}6 \\
        5{,}1 & 3{,}7 & -1{,}4 & | & -0{,}14
    \end{pmatrix}.
\end{equation}

Заданная точность для итерационного процесса: $\varepsilon = 10^{-5}$.

\section{Предварительный анализ числовых параметров системы}

Перед выполнением численных расчетов проведен аналитический анализ матрицы системы:
\begin{enumerate}
    \item \textbf{Вычисление определителя}:
    \begin{multline}
        \det(A) = 2{,}7 \cdot \big((-2{,}8) \cdot (-1{,}4) - 6{,}7 \cdot 3{,}7\big) - \\
        - 0{,}9 \cdot \big(4{,}5 \cdot (-1{,}4) - 6{,}7 \cdot 5{,}1\big) - \\
        - 1{,}5 \cdot \big(4{,}5 \cdot 3{,}7 - (-2{,}8) \cdot 5{,}1\big) = -66{,}321.
    \end{multline}
    Поскольку $\det(A) = -66{,}321 \neq 0$, матрица невырождена, следовательно, согласно теории решения СЛАУ, система совместна и имеет единственное решение;
    \item \textbf{Проверка диагонального преобладания}:
    \begin{itemize}
        \item строка 1: $|a_{11}| = 2{,}7$, сумма недиагональных элементов $|a_{12}| + |a_{13}| = 0{,}9 + 1{,}5 = 2{,}4$. Условие $2{,}7 > 2{,}4$ выполняется;
        \item строка 2: $|a_{22}| = 2{,}8$, сумма недиагональных элементов $|a_{21}| + |a_{23}| = 4{,}5 + 6{,}7 = 11{,}2$. Условие $2{,}8 > 11{,}2$ нарушено;
        \item строка 3: $|a_{33}| = 1{,}4$, сумма недиагональных элементов $|a_{31}| + |a_{32}| = 5{,}1 + 3{,}7 = 8{,}8$. Условие $1{,}4 > 8{,}8$ нарушено.
    \end{itemize}
\end{enumerate}

Таким образом, исходная система~\eqref{eq:var11} \textbf{не обладает свойством диагонального преобладания}. Непосредственное построение итерационной матрицы без преобразования приведет к норме $\|B\| > 1$ и расходимости итерационного процесса. Для выполнения задания 2 необходимо эквивалентное линейное комбинирование уравнений.

Исходные числовые данные хранятся во внешнем текстовом файле \texttt{11.in} (листинг~\ref{lst:input}).

\begin{listing}[H]
    \textfile{../11.in}
    \caption{Содержимое входного файла 11.in}
    \label{lst:input}
\end{listing}
